% ============================================================
%  TAARA Q.0 — Complete Technical Internals Document
%  Audience: Strong CS student, knows ML, new to security/quantum
%  Compile: pdflatex -interaction=nonstopmode taara_internals.tex
%           (run twice for cross-references)
% ============================================================
\documentclass[11pt,a4paper]{article}

% ── Packages ─────────────────────────────────────────────────
\usepackage[T1]{fontenc}
\usepackage[utf8]{inputenc}
\usepackage{lmodern}
\usepackage[margin=2.5cm]{geometry}
\usepackage{amsmath,amssymb,amsthm}
\usepackage{mathtools}
\usepackage{bm}
\usepackage{hyperref}
\usepackage{xcolor}
\usepackage{listings}
\usepackage{booktabs}
\usepackage{array}
\usepackage{tabularx}
\usepackage{multirow}
\usepackage{graphicx}
\usepackage{float}
\usepackage{enumitem}
\usepackage{tcolorbox}
\usepackage{tikz}
\usepackage{pgfplots}
\usepackage{caption}
\usepackage{subcaption}
\usepackage{fancyhdr}
\usepackage{titlesec}
\usepackage{parskip}
\usepackage{microtype}

\usetikzlibrary{quantikz2,arrows.meta,positioning,shapes,fit,decorations.pathmorphing,
                calc,matrix,backgrounds}
\pgfplotsset{compat=1.18}
\tcbuselibrary{skins,breakable}

% ── Colors ───────────────────────────────────────────────────
\definecolor{taaraRed}{HTML}{E94560}
\definecolor{taaraBlue}{HTML}{0F3460}
\definecolor{taaraGreen}{HTML}{00CC88}
\definecolor{taaraYellow}{HTML}{FFB347}
\definecolor{taaraGray}{HTML}{A0A0B0}
\definecolor{codebg}{HTML}{F5F5F5}
\definecolor{alertbg}{HTML}{FFF3CD}
\definecolor{notebg}{HTML}{E8F4FD}
\definecolor{mathbg}{HTML}{F0F8FF}
\definecolor{critbg}{HTML}{FDE8E8}

% ── Hyperref ─────────────────────────────────────────────────
\hypersetup{
  colorlinks=true, linkcolor=taaraBlue,
  citecolor=taaraRed, urlcolor=taaraBlue,
  pdftitle={TAARA Q.0 Internals},
  pdfauthor={IEPD / GoodWinSun}
}

% ── Custom tcolorboxes ────────────────────────────────────────
\tcbset{
  mybox/.style={
    enhanced, breakable, colback=#1!5!white, colframe=#1,
    fonttitle=\bfseries\small, left=6pt, right=6pt,
    top=4pt, bottom=4pt, boxrule=1pt
  }
}
\newtcolorbox{conceptbox}[2][]{
  mybox=taaraBlue, title={#2}, #1
}
\newtcolorbox{mathbox}[2][]{
  mybox=taaraRed, title={#2}, #1
}
\newtcolorbox{examplebox}[2][]{
  mybox=taaraGreen, title={#2}, #1
}
\newtcolorbox{warningbox}[1][]{
  mybox=taaraYellow, title={Key Insight}, #1
}
\newtcolorbox{codebox}[2][]{
  enhanced, breakable, colback=codebg, colframe=taaraGray!60,
  fonttitle=\ttfamily\small\bfseries, title={#2}, #1,
  boxrule=0.5pt
}

% ── Code listing style ────────────────────────────────────────
\lstset{
  backgroundcolor=\color{codebg},
  basicstyle=\ttfamily\footnotesize,
  breaklines=true, frame=single, rulecolor=\color{taaraGray!40},
  keywordstyle=\color{taaraBlue}\bfseries,
  commentstyle=\color{taaraGreen!70!black},
  stringstyle=\color{taaraRed},
  numbers=left, numberstyle=\tiny\color{taaraGray},
  numbersep=5pt, xleftmargin=10pt, tabsize=2
}

% ── Headers ──────────────────────────────────────────────────
\pagestyle{fancy}
\fancyhf{}
\lhead{\color{taaraGray}\small TAARA Q.0 — Technical Internals}
\rhead{\color{taaraGray}\small \thepage}
\renewcommand{\headrulewidth}{0.4pt}

% ── Section style ─────────────────────────────────────────────
\titleformat{\section}{\Large\bfseries\color{taaraBlue}}{}{0em}
  {\thesection\quad}[\vspace{-0.3em}\rule{\linewidth}{1.5pt}]
\titleformat{\subsection}{\large\bfseries\color{taaraBlue!80}}{}{0em}{\thesubsection\quad}
\titleformat{\subsubsection}{\normalsize\bfseries\color{taaraBlue!60}}{}{0em}{\thesubsubsection\quad}

% ── Math operators ────────────────────────────────────────────
\DeclareMathOperator{\MSE}{MSE}
\newcommand{\Rn}{\mathbb{R}^n}
\newcommand{\ket}[1]{|#1\rangle}
\newcommand{\bra}[1]{\langle #1|}
\newcommand{\braket}[2]{\langle #1|#2\rangle}
\newcommand{\norm}[1]{\left\lVert#1\right\rVert}
\newcommand{\mat}[1]{\mathbf{#1}}

% ============================================================
\begin{document}

% ── Title page ────────────────────────────────────────────────
\begin{titlepage}
\centering
\vspace*{2cm}

\begin{tikzpicture}
  \node[draw=taaraRed, line width=2pt, rounded corners=10pt,
        fill=taaraBlue!8, minimum width=14cm, minimum height=4cm] (box) {};
  \node[font=\Huge\bfseries\color{taaraRed}] at (box.center) {TAARA Q.0};
  \node[font=\large\color{taaraBlue}, below=0.3cm of box.center] {Complete Technical Internals};
\end{tikzpicture}

\vspace{1cm}
{\Large\color{taaraGray} From First Principles to Running Code}

\vspace{0.5cm}
\rule{12cm}{0.5pt}

\vspace{0.8cm}
{\large
\textbf{Topics Covered:}\\[0.4cm]
\begin{tabular}{ll}
§1 & SSH Protocol and Why Attacks Work \\
§2 & Why Behavioral Detection Beats Signatures \\
§3 & Reconstruction-Based Novelty Detection (TAARA Core) \\
§4 & The Autoencoder: Architecture and Mathematics \\
§5 & The PennyLane Quantum Circuit: Step by Step \\
§6 & Post-Quantum Cryptography and Kyber768 \\
§7 & MITRE ATT\&CK Framework and Parameter Mapping \\
§8 & GraphRAG: From Flat RAG to Knowledge Graphs \\
§9 & The Three Hidden Parameters: Why They Are Undetectable \\
§10 & Measurement, PDF Generation, and the Full Pipeline
\end{tabular}
}

\vspace{1.5cm}
\rule{12cm}{0.5pt}

\vspace{0.5cm}
{\normalsize\color{taaraGray}
IEPD / GoodWinSun · TAARA Q.0\\
\textit{``Prevent Crash, Preserve Cash''}\\[0.3cm]
Audience: CS student with ML background, new to cybersecurity and quantum
}
\end{titlepage}

\tableofcontents
\newpage

% ============================================================
\section{SSH Protocol and How Attacks Work at the Protocol Level}
% ============================================================

\subsection{What SSH Actually Is}

SSH (Secure Shell) is a cryptographic network protocol that replaces
plaintext protocols like Telnet and rlogin. But calling it ``encrypted
Telnet'' misses the threat model. Let's go from the TCP handshake to
where attacks live.

\begin{conceptbox}{The SSH Handshake — Layer by Layer}
\textbf{Step 1: TCP connection.}
The client opens a TCP connection to the server's port (default 22).
At this point, nothing is encrypted. The server sends its \emph{banner}:
\texttt{SSH-2.0-OpenSSH\_8.9p1}. An attacker scanning the internet
reads this banner and knows exactly what server software version is running.

\textbf{Step 2: Algorithm negotiation.}
Client and server exchange \texttt{KEXINIT} messages listing their
supported key exchange (kex) algorithms, ciphers, MACs, and compression
schemes. They pick the first matching pair. If the server is misconfigured
to support \texttt{diffie-hellman-group1-sha1} (768-bit DH, broken since 2015),
a client can downgrade to it.

\textbf{Step 3: Key exchange (Diffie-Hellman or ECDH).}
The client and server agree on a \emph{session key} without ever transmitting it.
After this step, all traffic is encrypted using that session key.
The server's \emph{host key} authenticates the server's identity (prevents MITM).

\textbf{Step 4: Authentication.}
The client proves its identity — via password, public key, or certificate.
This is where most attacks land.

\textbf{Step 5: Channel multiplexing.}
SSH multiplexes an arbitrary number of ``channels'' (shell, file transfer,
port forwarding) over one connection.
\end{conceptbox}

\subsection{Where the Attack Surface Lives}

\begin{figure}[H]
\centering
\begin{tikzpicture}[every node/.style={font=\small}]
  % Client
  \node[draw=taaraBlue, fill=taaraBlue!10, rounded corners, minimum width=2.5cm,
        minimum height=1cm] (cli) at (0,0) {Client};
  % Server
  \node[draw=taaraBlue, fill=taaraBlue!10, rounded corners, minimum width=2.5cm,
        minimum height=1cm] (srv) at (10,0) {Server};
  % Layers
  \foreach \y/\label/\col in {
    -1.5/TCP Handshake/taaraGray,
    -3/Algorithm Negotiation/taaraYellow,
    -4.5/Key Exchange + Auth/taaraRed,
    -6/Channel Operations/taaraGreen
  } {
    \draw[->, >=Stealth, color=\col, thick] (cli.south) -- ++(0,\y-0.3)
      -- ++(10,0) node[right, color=\col] {\label};
  }
  % Attack arrows
  \node[color=taaraRed, right, font=\footnotesize] at (5,-1.5)
    {\textit{banner reveals version}};
  \node[color=taaraRed, right, font=\footnotesize] at (5,-3.3)
    {\textit{downgrade attacks if weak algs enabled}};
  \node[color=taaraRed, right, font=\footnotesize] at (5,-5)
    {\textit{brute force, credential stuffing}};
\end{tikzpicture}
\caption{SSH handshake layers and where attacks occur.}
\end{figure}

The three principal attack vectors TAARA monitors:

\begin{enumerate}
\item \textbf{Brute force / credential stuffing} (MITRE T1110).
  A server with \texttt{PasswordAuthentication yes} will accept password
  guesses indefinitely if \texttt{MaxAuthTries} is not set. Attackers run
  tools like Hydra or Medusa, trying passwords from breach databases.
  The key technical fact: each failed attempt appears in \texttt{/var/log/auth.log}
  as a line like:\\
  \texttt{Failed password for root from 1.2.3.4 port 47291 ssh2}

\item \textbf{Privilege escalation via root login} (MITRE T1078).
  If \texttt{PermitRootLogin yes}, a successful authentication gives
  the attacker UID 0 immediately — no lateral movement required.

\item \textbf{Lateral movement via misconfigured forwarding} (MITRE T1021.004).
  SSH agent forwarding (\texttt{AllowAgentForwarding yes}) lets the
  attacker on a compromised intermediate host use the victim's SSH keys
  to pivot to other servers without ever seeing the private key material.
\end{enumerate}

\begin{examplebox}{Worked Example: What Happens in a Brute-Force Attack}
\textbf{Setup:} A server runs OpenSSH with \texttt{PasswordAuthentication yes}
and \texttt{MaxAuthTries 6} (the default).

\textbf{What the attacker's tool does:}
\begin{enumerate}
  \item Connects TCP to port 22. Reads banner: knows it's OpenSSH 8.9.
  \item Sends \texttt{SSH\_MSG\_USERAUTH\_REQUEST} for user \texttt{root},
        auth-type \texttt{password}, value \texttt{admin123}.
  \item Server replies \texttt{SSH\_MSG\_USERAUTH\_FAILURE}.
  \item Repeats up to 6 times per connection before being disconnected.
  \item Opens a new connection and tries 6 more passwords.
\end{enumerate}

\textbf{What TAARA records:} For each authentication window (say, 5 minutes),
TAARA extracts the feature vector $\mathbf{x}_t \in \mathbb{R}^7$ from auth.log:
\[
\mathbf{x}_t = [\underbrace{247}_{\text{failed}},\; \underbrace{2}_{\text{accepted}},\;
\underbrace{189}_{\text{invalid users}},\; \underbrace{1}_{\text{established}},\;
\underbrace{43}_{\text{unique src IPs}},\; \underbrace{5}_{\text{findings}},\;
\underbrace{125}_{\text{severity score}}]
\]

In a normal window, \texttt{failed\_logins}$\approx 3$, \texttt{unique\_src\_ips}$\approx 1$.
The residual $\Delta_t = \mathbf{x}_t - \hat{\mathbf{x}}_t$ is
far outside the memory basis — TAARA flags it as novel.
\end{examplebox}

% ============================================================
\section{Why Behavioral Detection Beats Signatures}
% ============================================================

\subsection{The Signature Approach and Its Fundamental Limit}

A signature-based IDS (like Snort or Suricata) contains rules of the form:
\begin{center}
  \texttt{alert tcp any any -> \$HOME\_NET 22 (msg:"SSH Brute Force"; \ldots)}
\end{center}

The rule fires if it sees a pattern it already knows. This creates a
fundamental asymmetry: \emph{the defender must know the attack before the attack
happens.} The attacker needs to be novel exactly once.

\begin{warningbox}
The signature approach asks: \textit{``Does this look like a known attack?''}\\
The behavioral approach asks: \textit{``Does this look like this specific system behaving normally?''}\\[0.3em]
These are different questions. The second one can catch attacks that have never been seen before.
\end{warningbox}

\subsection{The Threshold Problem}

Classical anomaly detectors (z-score, isolation forest, OCSVM) set a threshold $\tau$:
\[
\text{alert if } d(\mathbf{x}_t, \mu) > \tau
\]
where $\mu$ is the mean of observed behavior and $d$ is some distance function.

\textbf{The problem:} $\tau$ is a hyperparameter. If $\tau$ is too low,
you get false positives (alert fatigue). Too high, you miss attacks.
Worse: the optimal $\tau$ for one identity (a developer machine that runs
many processes) is completely different from another (a production database
server that barely changes). You end up tuning $\tau$ per-identity,
which requires labeled attack data you rarely have.

TAARA eliminates $\tau$ by replacing the fixed threshold with a
\emph{relative criterion} based on the identity's own history.

% ============================================================
\section{Reconstruction-Based Novelty Detection: The TAARA Core}
% ============================================================

\subsection{The Core Idea}

Imagine you have observed a system for 30 days. You have 30 ``behavioral
state'' vectors. Now on day 31, a new state arrives. TAARA asks:

\begin{center}
\textit{Can this new state be expressed as a linear combination of my prior observations?}
\end{center}

If yes: the new state is in the \emph{span} of prior behavior. Not novel.
If no: the new state contains a component orthogonal to everything seen before.
That orthogonal component is the residual. The larger it is (relative to
prior residuals), the more novel the behavior.

\subsection{Per-Identity Memory Basis}

For each monitored entity (a server, a user account, a container),
TAARA maintains a \emph{memory basis}:
\[
\mathcal{M}_u = \{\mathbf{m}_1, \mathbf{m}_2, \ldots, \mathbf{m}_k\}
\quad \mathbf{m}_i \in \mathbb{R}^n
\]
where $u$ is the identity label (e.g., \texttt{ssh\_server}, \texttt{user:alice}).

The basis matrix is:
\[
\mat{M} = [\mathbf{m}_1 \mid \mathbf{m}_2 \mid \cdots \mid \mathbf{m}_k]
\in \mathbb{R}^{n \times k}
\]
Columns of $\mat{M}$ are the stored behavioral states.
The \emph{column space} of $\mat{M}$, written $\text{col}(\mat{M})$, is the
subspace spanned by all previously observed behaviors.

\subsection{Least-Squares Projection}

Given a new observation $\mathbf{x}_t \in \mathbb{R}^n$, we find
the best approximation $\hat{\mathbf{x}}_t$ that lies in $\text{col}(\mat{M})$:

\begin{mathbox}{Equations 3–5: Reconstruction via Least-Squares Projection}
\[
\hat{\mathbf{x}}_t = \mat{M} \left(\mat{M}^\top \mat{M}\right)^{-1} \mat{M}^\top \mathbf{x}_t
\tag{3}
\]
\[
\boldsymbol{\Delta}_t = \mathbf{x}_t - \hat{\mathbf{x}}_t \tag{4}
\]
\[
\norm{\boldsymbol{\Delta}_t} = \sqrt{\sum_{i=1}^{n} \Delta_{t,i}^2} \tag{5}
\]
\end{mathbox}

\textbf{Why this formula?} The expression $\mat{M}(\mat{M}^\top\mat{M})^{-1}\mat{M}^\top$
is the \emph{orthogonal projection matrix} onto $\text{col}(\mat{M})$.
It's the unique linear operator that:
\begin{itemize}
  \item Is idempotent: $\mat{P}^2 = \mat{P}$ (projecting twice = projecting once)
  \item Is symmetric: $\mat{P}^\top = \mat{P}$
  \item Maps anything in $\text{col}(\mat{M})$ to itself
  \item Maps $\hat{\mathbf{x}}_t$ to the closest point in $\text{col}(\mat{M})$ to $\mathbf{x}_t$
\end{itemize}

\begin{examplebox}{Worked Numerical Example: Projection}
\textbf{Setup:} 3-dimensional behavior space ($n=3$).
Memory basis has 2 stored states:
\[
\mat{M} = \begin{bmatrix} 1 & 0 \\ 0 & 1 \\ 0 & 0 \end{bmatrix}
\]
These span the $xy$-plane. A new observation:
$\mathbf{x}_t = [2,\; 3,\; 5]^\top$

\textbf{Compute projection:}
\[
\mat{M}^\top\mat{M} = \begin{bmatrix} 1 & 0 \\ 0 & 1 \end{bmatrix}, \quad
(\mat{M}^\top\mat{M})^{-1} = \begin{bmatrix} 1 & 0 \\ 0 & 1 \end{bmatrix}
\]
\[
\hat{\mathbf{x}}_t = \mat{M} \cdot \mat{I} \cdot \mat{M}^\top \mathbf{x}_t
= \begin{bmatrix} 1 & 0 \\ 0 & 1 \\ 0 & 0 \end{bmatrix}
  \begin{bmatrix} 2 \\ 3 \end{bmatrix} = \begin{bmatrix} 2 \\ 3 \\ 0 \end{bmatrix}
\]
\textbf{Residual:}
$\boldsymbol{\Delta}_t = [2,3,5]^\top - [2,3,0]^\top = [0,0,5]^\top$

$\norm{\boldsymbol{\Delta}_t} = 5$. \textbf{Interpretation:} the new behavior
has a component of magnitude 5 in the $z$-direction, which the memory basis
has never seen. That ``$z$-axis'' behavior is completely new.
\end{examplebox}

\subsection{Threshold-Free Novelty Criterion}

\begin{mathbox}{Equation 6: Novelty Without a Fixed Threshold}
\[
\text{Novel}(t) = \mathbf{1}\left[\norm{\boldsymbol{\Delta}_t} > \max_{i < t} \norm{\boldsymbol{\Delta}_i}\right]
\tag{6}
\]
\end{mathbox}

This says: ``Flag this observation as novel if and only if its residual is
strictly larger than every residual seen before from this identity.''

\textbf{Why this is powerful:} There is no global threshold. A developer
machine that runs 400 processes may have a large normal residual norm of 15.0.
A production database that barely changes may have a normal residual norm of 0.3.
An attacker who causes a residual of 1.0 on the database is flagged.
The same attacker on the developer machine is not. This is correct:
the database doing something unusual is the signal; the developer machine
fluctuating is noise.

\subsection{Bootstrap Phase}

For a new identity (first connection), TAARA doesn't have a baseline yet.
The first \textbf{3 observations} are used to build the initial basis without
triggering alerts. After 3 observations, detection begins.

\begin{figure}[H]
\centering
\begin{tikzpicture}[
  node distance=1.2cm,
  box/.style={draw, rounded corners, fill=taaraBlue!10, minimum width=3.5cm,
               minimum height=0.7cm, font=\small}
]
  \node[box] (obs) {New observation $\mathbf{x}_t$};
  \node[box, below of=obs] (boot) {Bootstrap phase? ($k < 3$)};
  \node[box, below left=0.8cm and 0.5cm of boot, fill=taaraGreen!20] (add)
    {Add to basis, no alert};
  \node[box, below right=0.8cm and 0.5cm of boot, fill=taaraRed!10] (proj)
    {Project onto $\text{col}(\mat{M})$, compute $\boldsymbol{\Delta}_t$};
  \node[box, below of=proj] (thresh) {$\norm{\boldsymbol{\Delta}_t} > \max_{i<t}\norm{\boldsymbol{\Delta}_i}$?};
  \node[box, below left=0.8cm and 0.5cm of thresh, fill=taaraGreen!20] (notnovel)
    {Update basis, update max};
  \node[box, below right=0.8cm and 0.5cm of thresh, fill=taaraRed!20] (novel)
    {Novel! $\rightarrow$ Quantum validation};
  \draw[->, >=Stealth] (obs) -- (boot);
  \draw[->, >=Stealth] (boot) -- node[left, font=\scriptsize]{Yes} (add);
  \draw[->, >=Stealth] (boot) -- node[right, font=\scriptsize]{No} (proj);
  \draw[->, >=Stealth] (proj) -- (thresh);
  \draw[->, >=Stealth] (thresh) -- node[left, font=\scriptsize]{No} (notnovel);
  \draw[->, >=Stealth] (thresh) -- node[right, font=\scriptsize]{Yes} (novel);
\end{tikzpicture}
\caption{TAARA core detection pipeline.}
\end{figure}

% ============================================================
\section{The Autoencoder: Architecture and Mathematics}
% ============================================================

\subsection{Why Use an Autoencoder Here?}

An autoencoder does the same thing as the reconstruction above, but with a
\emph{non-linear} mapping learned from data. The key advantage: a 19-dimensional
raw feature vector may have a highly non-linear manifold of ``normal'' behavior.
A linear projection (§3) can miss non-linear deviations. The autoencoder learns
the shape of normal behavior non-linearly.

\subsection{Architecture}

The TAARA autoencoder (in \texttt{dna\_autoencoder.py}) has:

\begin{table}[H]
\centering
\begin{tabular}{@{}lllr@{}}
\toprule
Layer & Type & Activation & Dimension \\
\midrule
Input & --- & --- & 19 \\
Encoder hidden & Linear + Dropout(0.1) & ReLU & 64 \\
Bottleneck & Linear & ReLU & 64 \\
Decoder hidden & Linear + Dropout(0.1) & ReLU & 64 \\
Output & Linear & (none) & 19 \\
\bottomrule
\end{tabular}
\caption{TAARA autoencoder layer sizes.}
\end{table}

\begin{figure}[H]
\centering
\begin{tikzpicture}[scale=0.9, every node/.style={font=\small}]
  % Input neurons
  \foreach \i in {1,...,5} {
    \node[draw=taaraBlue, circle, minimum size=0.4cm, fill=taaraBlue!20]
      (x\i) at (0, {-\i*0.7+2.1}) {};
  }
  \node at (0, -2.5) {$\vdots$};
  \node[font=\footnotesize, below] at (0,-3.1) {19 inputs};

  % Encoder hidden
  \foreach \i in {1,...,6} {
    \node[draw=taaraYellow!80!black, circle, minimum size=0.4cm, fill=taaraYellow!30]
      (h1\i) at (2.5, {-\i*0.55+1.9}) {};
  }
  \node at (2.5, -2.0) {$\vdots$};
  \node[font=\footnotesize, below] at (2.5,-2.5) {64 neurons};

  % Bottleneck
  \foreach \i in {1,...,4} {
    \node[draw=taaraRed, circle, minimum size=0.4cm, fill=taaraRed!20]
      (bn\i) at (5, {-\i*0.6+1.4}) {};
  }
  \node at (5,-1.8) {$\vdots$};
  \node[font=\footnotesize, below] at (5,-2.3) {64 bottleneck};

  % Decoder hidden
  \foreach \i in {1,...,6} {
    \node[draw=taaraYellow!80!black, circle, minimum size=0.4cm, fill=taaraYellow!30]
      (h2\i) at (7.5, {-\i*0.55+1.9}) {};
  }
  \node at (7.5,-2.0) {$\vdots$};
  \node[font=\footnotesize, below] at (7.5,-2.5) {64 neurons};

  % Output
  \foreach \i in {1,...,5} {
    \node[draw=taaraBlue, circle, minimum size=0.4cm, fill=taaraGreen!30]
      (o\i) at (10, {-\i*0.7+2.1}) {};
  }
  \node at (10,-2.5) {$\vdots$};
  \node[font=\footnotesize, below] at (10,-3.1) {19 outputs ($\hat{\mathbf{x}}$)};

  % Connections (sampled)
  \foreach \i in {1,...,5} {
    \foreach \j in {1,...,6} {
      \draw[->, >=Stealth, taaraBlue!20, line width=0.3pt] (x\i) -- (h1\j);
    }
  }
  \foreach \i in {1,...,6} {
    \foreach \j in {1,...,4} {
      \draw[->, >=Stealth, taaraYellow!40, line width=0.3pt] (h1\i) -- (bn\j);
    }
  }
  \foreach \i in {1,...,4} {
    \foreach \j in {1,...,6} {
      \draw[->, >=Stealth, taaraRed!30, line width=0.3pt] (bn\i) -- (h2\j);
    }
  }
  \foreach \i in {1,...,6} {
    \foreach \j in {1,...,5} {
      \draw[->, >=Stealth, taaraGreen!30, line width=0.3pt] (h2\i) -- (o\j);
    }
  }
  % Labels
  \node[above, font=\small\bfseries\color{taaraBlue}] at (0, 2.5) {Input};
  \node[above, font=\small\bfseries\color{taaraYellow!80!black}] at (2.5, 2.5) {Enc.};
  \node[above, font=\small\bfseries\color{taaraRed}] at (5, 2.0) {Bottleneck};
  \node[above, font=\small\bfseries\color{taaraYellow!80!black}] at (7.5, 2.5) {Dec.};
  \node[above, font=\small\bfseries\color{taaraGreen!60!black}] at (10, 2.5) {Output};
  % Residual arrow
  \draw[<->, >=Stealth, taaraRed, thick, dashed]
    (0,-3.5) -- node[below, font=\small] {$\boldsymbol{\Delta} = \mathbf{x} - \hat{\mathbf{x}}$} (10,-3.5);
\end{tikzpicture}
\caption{Autoencoder architecture. The bottleneck forces the network to learn
a compact representation. Reconstruction error on \emph{new} data = anomaly signal.}
\end{figure}

\subsection{Training: Only on Benign Data}

\begin{mathbox}{Loss Function and Training Objective}
The autoencoder minimizes Mean Squared Error (MSE) reconstruction loss
over the training set $\mathcal{D} = \{\mathbf{x}_1, \ldots, \mathbf{x}_N\}$
of \textbf{benign-only} observations:
\[
\mathcal{L}(\theta) = \frac{1}{N} \sum_{i=1}^{N} \norm{\mathbf{x}_i - f_\theta(\mathbf{x}_i)}^2
\]
where $f_\theta: \mathbb{R}^{19} \to \mathbb{R}^{19}$ is the autoencoder
(encoder $\circ$ decoder), parameterized by weights $\theta$.

Adam optimizer: $\theta_{t+1} = \theta_t - \alpha \frac{\hat{m}_t}{\sqrt{\hat{v}_t} + \epsilon}$

where $\hat{m}_t$ and $\hat{v}_t$ are bias-corrected first and second moment estimates.
Learning rate $\alpha = 0.001$, $\epsilon = 10^{-8}$.
\end{mathbox}

\textbf{Why train only on benign data?} This is the core of semi-supervised
anomaly detection. The autoencoder learns to reconstruct \emph{normal} behavior.
When an attack occurs, the input is outside the distribution the network
was trained on. The network has no way to reconstruct it well, so the
reconstruction error spikes. We never needed labeled attack examples.

\subsection{The 19 Input Features}

From \texttt{taaraware\_manager.py}, the TaaraWare agent on the customer server
collects:

\begin{table}[H]
\centering
\begin{tabular}{@{}rll@{}}
\toprule
Index & Feature Name & What it Measures \\
\midrule
0 & \texttt{proc\_spawn\_rate} & Processes with runtime $<$ 60s \\
1 & \texttt{proc\_short\_lived\_ratio} & Fraction with runtime $<$ 5s \\
2 & \texttt{proc\_uid\_diversity} & Distinct UIDs running processes \\
3 & \texttt{proc\_root\_ratio} & Fraction of processes as root \\
4 & \texttt{proc\_cmd\_entropy} & Shannon entropy of command names \\
5 & \texttt{net\_outbound\_conn\_rate} & Total outbound connections \\
6 & \texttt{net\_unique\_dst\_ips} & Distinct destination IPs \\
7 & \texttt{net\_unique\_dst\_ports} & Distinct destination ports \\
8 & \texttt{net\_port\_entropy} & Shannon entropy of destination ports \\
9 & \texttt{net\_failed\_conn\_ratio} & Fraction of failed connections \\
10 & \texttt{cpu\_usage} & CPU utilization (\%) \\
11 & \texttt{memory\_usage} & Memory utilization (\%) \\
12 & \texttt{disk\_usage} & Disk utilization (\%) \\
13 & \texttt{temporal\_rhythm\_deviation} & Hidden (§9) \\
14 & \texttt{causal\_chain\_novelty} & Hidden (§9) \\
15 & \texttt{concealment\_signal} & Hidden (§9) \\
16--18 & (reserved) & future use \\
\bottomrule
\end{tabular}
\caption{TAARA's 19 behavioral features. Features 13–15 are hidden (§9).}
\end{table}

\subsection{Shannon Entropy: Why It Matters}

Feature 4 and 8 use Shannon entropy $H$:
\[
H(X) = -\sum_{x \in \mathcal{X}} p(x) \log_2 p(x)
\]

\textbf{Intuition for security:} A normal server runs the same 15 processes
repeatedly: \texttt{sshd}, \texttt{nginx}, \texttt{postgres}, etc.
Low entropy. An attacker running a port scan launches dozens of ephemeral
probe processes with random names. High entropy. A cryptominer spawns
one long-running process with an unusual name — changes the entropy signature
without changing the count.

\begin{examplebox}{Entropy Example: Normal vs Under Attack}
\textbf{Normal (10 windows, 3 unique commands: sshd, nginx, postgres):}\\
$p(\texttt{sshd}) = 0.6,\; p(\texttt{nginx}) = 0.3,\; p(\texttt{postgres}) = 0.1$
\[
H = -(0.6\log_2 0.6 + 0.3\log_2 0.3 + 0.1\log_2 0.1) \approx 1.30 \text{ bits}
\]

\textbf{Under attack (port scan spawns 8 unique probe commands, uniform):}\\
$p(\text{each}) = 0.125$
\[
H = -8 \cdot 0.125 \log_2 0.125 = -8 \cdot 0.125 \cdot (-3) = 3.0 \text{ bits}
\]

The entropy feature jumps from 1.30 to 3.0 — a 2.3$\times$ increase.
The autoencoder, trained only on windows with entropy $\approx 1.30$,
will fail to reconstruct this window and produce high $\|\Delta\|$.
\end{examplebox}

% ============================================================
\section{The PennyLane Quantum Circuit: Step by Step}
% ============================================================

\subsection{What Quantum Computing Actually Is (Honest)}

A classical bit is either 0 or 1. A qubit is described by a state:
\[
\ket{\psi} = \alpha\ket{0} + \beta\ket{1}, \quad |\alpha|^2 + |\beta|^2 = 1
\]
where $\alpha, \beta \in \mathbb{C}$ are \emph{probability amplitudes}.
$|\alpha|^2$ is the probability of measuring 0; $|\beta|^2$ of measuring 1.

\textbf{Important for TAARA:} We run PennyLane on a \emph{classical simulator}.
There is no quantum speedup claim. The value is purely in the
\emph{mathematical structure} of quantum fidelity as a principled,
parameter-free directionality measure.

\subsection{4-Qubit System: The State Space}

With 4 qubits, the computational basis has $2^4 = 16$ states:
$\ket{0000}, \ket{0001}, \ldots, \ket{1111}$.

A general 4-qubit state is:
\[
\ket{\psi} = \sum_{i=0}^{15} c_i \ket{i}, \qquad \sum_{i=0}^{15} |c_i|^2 = 1
\]

This 16-dimensional complex unit vector is exactly what we need to encode
our residual $\boldsymbol{\Delta}_t \in \mathbb{R}^n$.

\subsection{Amplitude Encoding}

To encode a real vector $\mathbf{v} \in \mathbb{R}^{16}$ into a quantum state:
\begin{enumerate}
  \item Normalize: $\hat{\mathbf{v}} = \mathbf{v} / \norm{\mathbf{v}}$
  \item Pad to 16 dimensions if needed (residual has $n \leq 7$ dims in TAARA)
  \item Set: $c_i = \hat{v}_i$ (real amplitudes)
\end{enumerate}

The resulting state $\ket{\psi} = \sum_{i=0}^{15} \hat{v}_i \ket{i}$
encodes the \emph{direction} of $\mathbf{v}$ in Hilbert space.
The magnitude is lost (everything is normalized to the unit sphere).
This is intentional: we care about \emph{which direction} the residual points,
not how large it is.

\subsection{The Circuit}

The TAARA quantum circuit (from \texttt{quantum\_engine.py}):

\begin{figure}[H]
\centering
\begin{tikzpicture}
  % Define qubit lines
  \foreach \q/\label in {0/q_3, 1/q_2, 2/q_1, 3/q_0} {
    \node[left, font=\small] at (-0.3, {-\q*1.2}) {$\ket{\label}$};
    \draw[thick] (0, {-\q*1.2}) -- (12, {-\q*1.2});
  }
  % AmplitudeEmbedding (gate box)
  \node[draw=taaraBlue, fill=taaraBlue!15, minimum width=1.2cm, minimum height=4.8cm,
        rounded corners=3pt, font=\scriptsize\bfseries] at (1, -1.8) {$\ket{\psi}$\\AmpEmb};

  % Hadamard gates
  \foreach \q in {0,1,2,3} {
    \node[draw=taaraYellow!80!black, fill=taaraYellow!30, minimum size=0.5cm,
          font=\scriptsize\bfseries] at (2.8, {-\q*1.2}) {H};
  }

  % Ring CNOT: 0->1, 1->2, 2->3, 3->0
  % CNOT 0->1
  \filldraw[taaraRed] (4.5, 0) circle (0.12cm);
  \draw[taaraRed, thick] (4.5, 0) -- (4.5, -1.2);
  \draw[taaraRed, thick] (4.5,-1.2) circle (0.2cm);
  \draw[taaraRed, thick] (4.3,-1.2) -- (4.7,-1.2);
  % CNOT 1->2
  \filldraw[taaraRed] (5.5, -1.2) circle (0.12cm);
  \draw[taaraRed, thick] (5.5,-1.2) -- (5.5,-2.4);
  \draw[taaraRed, thick] (5.5,-2.4) circle (0.2cm);
  \draw[taaraRed, thick] (5.3,-2.4) -- (5.7,-2.4);
  % CNOT 2->3
  \filldraw[taaraRed] (6.5, -2.4) circle (0.12cm);
  \draw[taaraRed, thick] (6.5,-2.4) -- (6.5,-3.6);
  \draw[taaraRed, thick] (6.5,-3.6) circle (0.2cm);
  \draw[taaraRed, thick] (6.3,-3.6) -- (6.7,-3.6);
  % CNOT 3->0 (ring)
  \filldraw[taaraRed] (7.5, -3.6) circle (0.12cm);
  \draw[taaraRed, thick] (7.5,-3.6) -- (7.5,0.5) -- (7.8,0.5) -- (7.8,0);
  \draw[taaraRed, thick] (7.8,0) circle (0.2cm);
  \draw[taaraRed, thick] (7.6,0) -- (8.0,0);

  % RX RY RZ on each qubit
  \foreach \q in {0,1,2,3} {
    \node[draw=taaraGreen!60!black, fill=taaraGreen!20, minimum width=0.45cm,
          minimum height=0.35cm, font=\tiny\bfseries]
      at (9.3, {-\q*1.2}) {RX};
    \node[draw=taaraGreen!60!black, fill=taaraGreen!20, minimum width=0.45cm,
          minimum height=0.35cm, font=\tiny\bfseries]
      at (10.1, {-\q*1.2}) {RY};
    \node[draw=taaraGreen!60!black, fill=taaraGreen!20, minimum width=0.45cm,
          minimum height=0.35cm, font=\tiny\bfseries]
      at (10.9, {-\q*1.2}) {RZ};
  }
  % Measurement symbols
  \foreach \q in {0,1,2,3} {
    \node[draw, fill=white, minimum size=0.4cm, font=\tiny] at (11.9, {-\q*1.2}) {M};
  }
  % Labels below
  \node[font=\footnotesize, below] at (1,-4.2) {Amp. Embed.};
  \node[font=\footnotesize, below] at (2.8,-4.2) {Hadamard};
  \node[font=\footnotesize, below, text=taaraRed] at (6,-4.2) {Ring CNOT};
  \node[font=\footnotesize, below] at (10.1,-4.2) {Rotations ($\pi/4$)};
  \node[font=\footnotesize, below] at (11.9,-4.2) {Measure};
\end{tikzpicture}
\caption{TAARA 4-qubit circuit. Depth: 5 gate layers. Parameters: 12 fixed rotation angles.}
\end{figure}

\subsubsection{Gate 1: AmplitudeEmbedding}

Encodes the normalized residual vector into the initial state:
\[
\text{AmpEmb}(\hat{\mathbf{v}}): \ket{0000} \mapsto \sum_{i=0}^{15} \hat{v}_i \ket{i}
\]
This is implemented by PennyLane as a sequence of single-qubit and
two-qubit rotations that prepare the exact target amplitude distribution.

\subsubsection{Gate 2: Hadamard}

The Hadamard gate is:
\[
H = \frac{1}{\sqrt{2}}\begin{pmatrix} 1 & 1 \\ 1 & -1 \end{pmatrix}
\]
Applied to each qubit, it creates superposition. More importantly for fidelity:
it spreads the amplitude information across all basis states, ensuring that
small differences in the input vector produce measurable differences
in the output state.

\subsubsection{Gate 3: Ring CNOT (Entanglement)}

The CNOT (controlled-NOT) gate:
\[
\text{CNOT} = \begin{pmatrix} 1 & 0 & 0 & 0 \\ 0 & 1 & 0 & 0 \\ 0 & 0 & 0 & 1 \\ 0 & 0 & 1 & 0 \end{pmatrix}
\]
In the ring topology (qubit 0 controls 1, 1 controls 2, 2 controls 3, 3 controls 0),
this creates \emph{entanglement} between all qubits. The result: the quantum state
now encodes correlations between features. An attacker who is trying to manipulate
one dimension of the feature vector must simultaneously manipulate all 4 qubits
(which requires knowing the full Kyber768 shared secret — see §6).

\subsubsection{Gates 4–6: Parameterized Rotations}

RX, RY, RZ with angle $\theta = \pi/4 \approx 0.785$ radians:
\[
R_X(\theta) = \begin{pmatrix} \cos(\theta/2) & -i\sin(\theta/2) \\ -i\sin(\theta/2) & \cos(\theta/2) \end{pmatrix}
= \begin{pmatrix} 0.924 & -0.383i \\ -0.383i & 0.924 \end{pmatrix}
\]
These rotations spread amplitude between $\ket{0}$ and $\ket{1}$ components
of each qubit. The angle $\pi/4$ is chosen as the geometric midpoint of
the rotation space — no fine-tuning, no free parameters.

\subsection{Quantum Fidelity: The Core Measurement}

\begin{mathbox}{Equations 7–9: Fidelity in TAARA}
Encode residual $\boldsymbol{\Delta}_t$ into quantum state:
\[
\hat{\boldsymbol{\Delta}} = \boldsymbol{\Delta}_t / \norm{\boldsymbol{\Delta}_t}, \quad
\ket{\psi_t} = \sum_{i=0}^{15} \hat{\Delta}_i \ket{i} \tag{7-8}
\]
Fidelity between current state and a memory state:
\[
F(\ket{\psi_t}, \ket{\psi_m}) = |\braket{\psi_t}{\psi_m}|^2 \tag{9}
\]
Minimum fidelity over all memory states:
\[
F_{\min}(t) = \min_{m \in \mathcal{M}_u} F(\ket{\psi_t}, \ket{\psi_m})
\]
\textbf{Quantum-confirmed novelty:} $F_{\min}(t) < 0.5$
\end{mathbox}

\textbf{Why 0.5?} This threshold is geometric, not empirical.
$F = 1.0$ means the two states are identical (same direction).
$F = 0.0$ means the states are orthogonal (completely different directions).
$F = 0.5$ is the geometric midpoint: the current state is as ``close'' to
the memory as it is ``far.'' Below 0.5, the new state is more orthogonal
than parallel — it genuinely occupies a new direction in behavior space.

\begin{examplebox}{Fidelity Worked Example}
\textbf{Normal day:} Residual $\boldsymbol{\Delta}_t = [1, 0, 0, 0]^\top$,
normalized: $\hat{\boldsymbol{\Delta}} = [1, 0, 0, 0]^\top$.
Memory state $\ket{\psi_m}$ was encoded from $[0.9, 0.4, 0, 0]^\top$, normalized: $\approx [0.914, 0.406, 0, 0]^\top$.

Fidelity $F = |\braket{\psi_t}{\psi_m}|^2 = |1 \cdot 0.914 + 0 + 0 + 0|^2 = 0.914^2 = 0.835$.

$F = 0.835 > 0.5$: normal, within baseline direction.

\textbf{Attack day:} Residual $\boldsymbol{\Delta}_t = [0, 0, 1, 0]^\top$ (novelty in 3rd dimension).
Fidelity $F = |0 \cdot 0.914 + 0 \cdot 0.406 + 1 \cdot 0 + 0|^2 = 0.0$.

$F = 0.0 < 0.5$: quantum-confirmed novel. The new residual is completely
orthogonal to every prior memory state.
\end{examplebox}

\subsection{Implementation: How the Fidelity Circuit Works in \texttt{query\_knowledge\_base.py}}

For the GraphRAG use case (scoring config deviation from policy), the circuit
is slightly different. It computes the probability of measuring $\ket{0000}$
after applying one amplitude embedding followed by the adjoint of the other:

\[
F = |\braket{\psi_a}{\psi_b}|^2 = \left|\bra{0000}\, U_b^\dagger U_a \,\ket{0000}\right|^2
\]

where $U_a$ encodes vector $\mathbf{a}$ and $U_b^\dagger$ is the inverse
(adjoint) of the encoding for $\mathbf{b}$. The $\ket{0000}$ outcome probability
from the composed circuit equals the inner product squared.
This uses PennyLane's \texttt{qml.adjoint()} to efficiently compute the overlap.

% ============================================================
\section{Post-Quantum Cryptography and Kyber768}
% ============================================================

\subsection{Why Classical Cryptography Is Threatened}

Classical public-key cryptography (RSA, ECDH) relies on computational problems
that are hard for classical computers but \emph{polynomial} for quantum computers
running Shor's algorithm:
\begin{itemize}
  \item RSA: integer factorization. Classical: sub-exponential. Shor: $O((\log N)^3)$.
  \item ECDH: discrete logarithm on elliptic curves. Same vulnerability.
\end{itemize}

This means that a sufficiently powerful quantum computer could break
all existing RSA and ECC keys. The threat is already present via
``harvest now, decrypt later'' attacks: adversaries are storing encrypted
traffic today to decrypt it once quantum computers arrive.

\subsection{Lattice-Based Cryptography and ML-KEM}

Kyber768 (officially \emph{ML-KEM} under NIST FIPS 203) is a \emph{Key
Encapsulation Mechanism} based on the hardness of the
\textbf{Module Learning With Errors (M-LWE)} problem.

\subsubsection{The Learning With Errors Problem}

LWE: given many pairs $(\mathbf{a}_i, b_i = \langle \mathbf{a}_i, \mathbf{s} \rangle + e_i)$
where $\mathbf{a}_i \in \mathbb{Z}_q^n$ is random, $\mathbf{s} \in \mathbb{Z}_q^n$ is a secret,
and $e_i$ is small noise, find $\mathbf{s}$.

\textbf{Why it's hard:} The noise corrupts the linear system. Without noise,
you'd just solve $n$ equations for $\mathbf{s}$. With noise, the best
known algorithms (lattice reduction: BKZ, LLL) take exponential time
in the lattice dimension. No quantum speedup beyond a quadratic Grover speedup
is known.

\begin{mathbox}{LWE in One Line}
\[
b_i = \underbrace{\langle \mathbf{a}_i, \mathbf{s} \rangle}_{\text{linear}} + e_i \pmod{q}
\]
Finding $\mathbf{s}$ from $\{(\mathbf{a}_i, b_i)\}$ is the LWE problem.
Changing $q$, $n$, or the noise distribution $\chi$ adjusts security level.
\end{mathbox}

\subsubsection{Module LWE and Kyber768 Parameters}

Kyber768 uses a \emph{module} structure: instead of vectors over $\mathbb{Z}_q$,
it uses polynomials in $\mathbb{Z}_q[x]/(x^{256}+1)$.
This gives better key sizes while maintaining LWE hardness.

\begin{table}[H]
\centering
\begin{tabular}{@{}llr@{}}
\toprule
Parameter & Value & Notes \\
\midrule
Module rank $k$ & 3 & (Kyber512: 2, Kyber1024: 4) \\
Polynomial degree $n$ & 256 & ring $\mathbb{Z}_q[x]/(x^{256}+1)$ \\
Modulus $q$ & 3329 & prime \\
Noise parameter $\eta$ & 2 & centered binomial distribution \\
Public key size & 1184 bytes & \\
Secret key size & 2400 bytes & \\
Ciphertext size & 1088 bytes & \\
Shared secret size & 32 bytes & \\
Classical security & $\geq 161$ bits & \\
Quantum security & $\geq 153$ qubits & \\
\bottomrule
\end{tabular}
\caption{Kyber768 parameters (NIST FIPS 203 Level 3).}
\end{table}

\subsection{How TAARA Uses Kyber768}

When TaaraWare is deployed to a customer server, the TAARA Command Center runs:

\begin{codebox}{Key Generation (taaraware\_manager.py:587--607)}
\begin{lstlisting}[language=Python]
import oqs  # liboqs-python: NIST-certified implementation
kem = oqs.KeyEncapsulation('Kyber768')
public_key = kem.generate_keypair()  # generates (pk, sk) internally
ciphertext, shared_secret = kem.encap_secret(public_key)
# shared_secret is 32 bytes — exists only in RAM, never written to disk
\end{lstlisting}
\end{codebox}

\begin{enumerate}
  \item Command Center generates a Kyber768 key pair $(pk, sk)$.
  \item \textbf{Encapsulation:} $(\text{ciphertext}, K) \leftarrow \text{Encaps}(pk)$.
        The 32-byte shared secret $K$ never leaves Command Center RAM.
  \item The \textbf{ciphertext} is sent to the deployed TaaraWare agent
        (or stored for future use). The \textbf{public key} and fingerprint
        are stored in \texttt{models/client\_keys.json}.
  \item For decapsulation (if agent had the private key):
        $K \leftarrow \text{Decaps}(sk, \text{ciphertext})$.
        Both sides now share $K$ without it ever being transmitted.
\end{enumerate}

\textbf{Current use in TAARA:} The shared secret $K$ is used as a seed
for the per-client feature transform (§9). Even if an attacker knows
the exact TAARA algorithm, they cannot compute what ``normal'' looks like
for a specific client without $K$.

% ============================================================
\section{MITRE ATT\&CK and Parameter Mapping}
% ============================================================

\subsection{What MITRE ATT\&CK Is}

MITRE ATT\&CK is a structured knowledge base of adversary tactics,
techniques, and procedures (TTPs) observed in real attacks.

\begin{table}[H]
\centering
\begin{tabular}{@{}lll@{}}
\toprule
Level & Example & Description \\
\midrule
Tactic & TA0006: Credential Access & The goal (what the attacker wants) \\
Technique & T1110: Brute Force & The method (how they achieve the goal) \\
Sub-technique & T1110.001: Password Guessing & Specific variant \\
Procedure & Hydra SSH brute-force & Specific tool / implementation \\
\bottomrule
\end{tabular}
\end{table}

\subsection{TAARA's MITRE Mapping}

From \texttt{taara\_analysis.py}, TAARA maps each detected finding to
its ATT\&CK entry automatically:

\begin{table}[H]
\centering
\footnotesize
\begin{tabular}{@{}p{4cm}lp{4.5cm}p{3cm}@{}}
\toprule
Configuration Issue & Tactic & Technique & Real-World Impact \\
\midrule
\texttt{PasswordAuthentication yes} & TA0006 Cred. Access & T1110 Brute Force & Unlimited password guessing \\
\texttt{PermitRootLogin yes} & TA0003 Persistence & T1078 Valid Accounts & Any login = full root = game over \\
Firewall disabled (ufw/iptables) & TA0005 Def. Evasion & T1562.004 Firewall Off & Entire internet can reach all ports \\
Audit/log missing & TA0005 Def. Evasion & T1562.002 Logs off & Attack is invisible, unrecoverable \\
\texttt{X11Forwarding yes} & TA0008 Lateral Mov. & T1021.004 SSH & GUI forwarding opens MITM path \\
Open port exposure & TA0043 Recon & T1046 Net. Discovery & Each port is an additional attack surface \\
Empty password & TA0006 Cred. Access & T1110.001 Pwd Guess & Guessed instantly, zero brute-force needed \\
\bottomrule
\end{tabular}
\caption{TAARA's ATT\&CK mappings for SSH and firewall findings.}
\end{table}

\textbf{Why this matters:} A finding like ``\texttt{PasswordAuthentication yes}''
by itself is just a configuration note. Mapping it to \textbf{T1110 + TA0006}
tells the security team: this enables Credential Access via Brute Force,
which is a tactic used in 73\% of all server breaches. The path from
misconfiguration to real-world attack is explicit.

% ============================================================
\section{GraphRAG: From Flat RAG to Knowledge Graphs}
% ============================================================

\subsection{The Problem with Flat RAG}

Standard RAG (Retrieval-Augmented Generation) works like this:
\begin{enumerate}
  \item Build an index of document chunks (embed each chunk as a vector).
  \item At query time: embed the query, find nearest chunks by cosine similarity.
  \item Feed retrieved chunks as context to the LLM.
\end{enumerate}

\textbf{Limitation:} Flat RAG retrieves by \emph{similarity}, not by
\emph{causality}. If your SSH config says \texttt{PermitRootLogin yes},
a flat RAG might retrieve the OWASP paragraph about ``use strong passwords.''
But what you really need to know is: this specific misconfiguration
\emph{enables} an exploit chain that \emph{leads to} full server compromise.
Flat RAG has no concept of ``leads to.''

\subsection{GraphRAG: Adding Causal Structure}

TAARA's knowledge graph is a directed graph $G = (V, E)$ where:
\begin{itemize}
  \item $V$ = security concepts (misconfigurations, attack techniques,
        consequences, mitigations)
  \item $E$ = typed edges: \texttt{enables}, \texttt{leads\_to},
        \texttt{mitigated\_by}, \texttt{exploits}
\end{itemize}

\begin{figure}[H]
\centering
\begin{tikzpicture}[
  every node/.style={font=\small},
  misconfig/.style={draw=taaraRed, fill=taaraRed!15, rounded corners, minimum width=3cm,
                    minimum height=0.6cm},
  attack/.style={draw=taaraYellow!80!black, fill=taaraYellow!20, rounded corners,
                 minimum width=3cm, minimum height=0.6cm},
  consequence/.style={draw=taaraRed!80, fill=taaraRed!30, rounded corners,
                      minimum width=3cm, minimum height=0.6cm},
  mitigation/.style={draw=taaraGreen!60!black, fill=taaraGreen!15, rounded corners,
                     minimum width=3cm, minimum height=0.6cm},
  edge label/.style={font=\tiny\itshape\color{taaraBlue}}
]
  \node[misconfig] (pw) at (0,0) {PasswordAuth yes};
  \node[attack] (bf) at (4,-1) {T1110 Brute Force};
  \node[attack] (va) at (4,1) {T1078 Valid Accounts};
  \node[consequence] (cc) at (8,0) {Full Server Compromise};
  \node[mitigation] (fix) at (0,-2) {Set PasswordAuth no};

  \draw[->, >=Stealth, taaraRed, thick] (pw) -- node[edge label, above]{enables} (bf);
  \draw[->, >=Stealth, taaraRed, thick] (pw) -- node[edge label, below]{enables} (va);
  \draw[->, >=Stealth, taaraRed, thick] (bf) -- node[edge label, above]{leads\_to} (cc);
  \draw[->, >=Stealth, taaraRed, thick] (va) -- node[edge label, below]{leads\_to} (cc);
  \draw[->, >=Stealth, taaraGreen!60!black, thick] (pw) -- node[edge label, right]{mitigated\_by} (fix);
\end{tikzpicture}
\caption{Knowledge graph fragment. A single misconfiguration enables two
attack paths, both leading to the same consequence. Flat RAG cannot
represent this causal structure.}
\end{figure}

\subsection{The Full TAARA GraphRAG Pipeline}

For the \emph{behavioral} scanner (platform scan / OHA):
\begin{enumerate}
  \item \textbf{Pattern scan:} Regex rules match known misconfigs in the config text.
  \item \textbf{Semantic search (FAISS):} Embed the config text using
        \texttt{all-MiniLM-L6-v2} (384-dim sentence vectors).
        FAISS approximate nearest neighbor search finds the 8 most similar
        policy chunks from 2,444 indexed chunks.
  \item \textbf{Quantum fidelity:} For each matched chunk,
        compute $F(\text{config embedding}, \text{policy embedding})$.
        Chunks with $F < 0.5$ are \emph{genuinely} in an unsafe direction.
  \item \textbf{Graph traversal:} For each matched node in the graph,
        traverse outgoing \texttt{enables}/\texttt{leads\_to} edges
        (up to depth 3) to build the propagation chain.
  \item \textbf{Mitigation lookup:} Follow \texttt{mitigated\_by} edges
        to retrieve fix descriptions.
\end{enumerate}

For the \emph{dependency} scanner (repo/code scan):
\begin{enumerate}
  \item \textbf{OSV query:} Get real CVE data for each dependency version.
  \item \textbf{FAISS retrieval:} Embed each CVE description,
        retrieve top-3 matching policy chunks.
  \item \textbf{LLM synthesis (Groq LLaMA-3.3-70B):}
        Build prompt from: CVE list + exploit chains + policy context.
        LLM generates a 3–4 sentence dependency risk analysis.
  \item \textbf{Quantum fidelity on the answer:}
        Embed the LLM's answer. Compare against the mean embedding
        of retrieved policy chunks.
        $F(\text{answer vec}, \text{policy mean vec})$ tells us:
        does the LLM's answer point in the same direction as established policy?
        $F > 0.6$: well-aligned. $F < 0.35$: review the answer.
\end{enumerate}

\begin{mathbox}{Why Fidelity on the LLM Answer?}
The LLM can hallucinate or drift from the security policy context
even when given correct facts. Instead of doing keyword matching
(``did the LLM mention NIST?''), TAARA embeds the answer and
computes how ``close'' it is to the policy space geometrically.
This is a principled, model-agnostic way to verify that the
LLM's guidance is aligned with the knowledge base — without a human reviewer.
\end{mathbox}

\subsection{FAISS: How Approximate Nearest Neighbor Works}

FAISS (Facebook AI Similarity Search) stores 2,444 policy vectors,
each in $\mathbb{R}^{384}$. Index type: Flat (exact L2 / inner product).

Given query vector $\mathbf{q} \in \mathbb{R}^{384}$, FAISS returns:
\[
\text{top-}k = \underset{i}{\arg\text{sort}} \left(-\mathbf{q}^\top \mathbf{p}_i\right)
\quad \text{(inner product = cosine similarity for normalized vectors)}
\]

With the \texttt{Flat} index and 2,444 vectors, this is brute-force
$O(384 \times 2444)$ floating point operations per query — about $10^6$ FLOPs,
fast on any CPU.

% ============================================================
\section{The Three Hidden Parameters: Why They Are Undetectable}
% ============================================================

\subsection{The Design Principle}

An attacker who reads the TAARA source code knows which 13 features
(indices 0–12) are collected. The attacker could craft behavior that
produces ``normal'' values for all 13 visible features while performing
a malicious action. This is the \emph{evasion problem}.

TAARA's response: add three additional features that measure
\emph{meta-properties of behavior history} — things that require
knowledge of the prior history that only the TAARA system has.

From \texttt{taaraware\_manager.py:252–373}:

\subsection{Feature 13: Temporal Rhythm Deviation}

\textbf{What it measures:} Whether the current hour-of-day has historically
been quiet on this server.

\textbf{Algorithm:}
\begin{enumerate}
  \item Load the last 100 feature buffer entries.
  \item Build a histogram: $H[h] = \text{count of entries with timestamp in hour } h$.
  \item Current hour fraction: $p_h = H[\text{current hour}] / 100$.
  \item Signal: $\texttt{temporal\_rhythm\_deviation} = 1.0$ if $p_h < 0.05$, else $0.0$.
\end{enumerate}

\textbf{Why it's hard to evade:}
The server's natural rhythm accumulates over months. If the server is
always quiet between 2–5am (because it's a production database used during
business hours), an attacker who stages an attack at 3am triggers a flag.
The attacker cannot know the server's exact hourly distribution without
having monitored it for months themselves. And if they \emph{have} been
monitoring it for months, they've been inside the network long enough that
other TAARA signals would have caught them.

\subsection{Feature 14: Causal Chain Novelty}

\textbf{What it measures:} Whether any parent→child process pair is new
relative to the last 50 observation windows.

\textbf{Algorithm:}
\begin{enumerate}
  \item Run \texttt{ps -eo ppid,pid,comm} to get all running processes.
  \item Hash each (parent\_comm, child\_comm) pair:
        $h = \text{MD5}(\texttt{ppid:comm})$.
  \item Collect all hashes seen in the last 50 buffer entries.
  \item $\texttt{causal\_chain\_novelty} = \frac{|\text{new pairs}|}{|\text{current pairs}|}$.
\end{enumerate}

\textbf{Why it's hard to evade:}
Every attack tool (a privilege escalation exploit, a reverse shell,
a cryptominer) has a characteristic process lineage. If \texttt{sshd}
spawns \texttt{bash} spawns \texttt{wget} — that's normal. If \texttt{sshd}
spawns \texttt{bash} spawns \texttt{nmap} — that's a new pair. The full
hash of (parent, child) combinations creates a unique fingerprint of
``what processes beget what.'' This is particularly effective against
living-off-the-land attacks that use legitimate OS tools: even if the
attacker uses \texttt{curl} (a legit tool), the parent→child chain
\texttt{sshd→bash→curl} may never have occurred on this server before.

\subsection{Feature 15: Concealment Signal}

This is a composite of three sub-signals that each detect different
forms of evidence destruction:

\begin{table}[H]
\centering
\begin{tabular}{@{}p{3.5cm}p{5.5cm}p{4cm}@{}}
\toprule
Sub-signal & Trigger & Attack this catches \\
\midrule
Auth log stagnation & auth.log file size didn't grow during business hours & Attacker deleting/truncating auth logs \\
Bash history shrinkage & \texttt{\~{}/.bash\_history} line count decreased & Attacker running \texttt{history -c} \\
Cron silence & Scheduled jobs existed but cron log not updated in 10 min & Ransomware disabling scheduled security jobs \\
\bottomrule
\end{tabular}
\caption{The three concealment sub-signals that compose Feature 15.}
\end{table}

$\texttt{concealment\_signal} = \frac{s_1 + s_2 + s_3}{3}$

\textbf{Why it's brilliant and why it's undetectable:}
An attacker covering their tracks will fire at least one of these.
If they delete auth logs ($s_1 = 1.0$) but leave bash history unchanged ($s_2 = 0.0$)
and don't touch cron ($s_3 = 0.0$), the signal is $1/3 = 0.33$ — a moderate anomaly.
If they do all three, signal = 1.0.

The reason this is hard to evade: to avoid triggering $s_1$, the attacker
must \emph{write fake auth log entries} to keep the file growing while they work.
But writing fake log entries creates a distinctive pattern in the log content
that would be caught by the policy deviation scanner (§8). To avoid $s_2$,
they must leave their commands in bash history — which exposes their tooling.
To avoid $s_3$, they must let cron jobs run — which may trigger other detection.
Each evasion creates a new detectable signal.

\subsection{The Kyber-Keyed Transform: Making Hidden Features Truly Unguessable}

Even if an attacker reads the source code and knows all three hidden features,
they still cannot compute ``normal'' for a specific client. The function
\texttt{apply\_client\_transform()} applies a per-client offset before
quantum encoding:

\begin{mathbox}{Client Transform: Lattice-Hard Evasion Barrier}
\[
\mathbf{v}_{\text{keyed}} = \frac{\mathbf{v} + \boldsymbol{\delta}_K}
  {\norm{\mathbf{v} + \boldsymbol{\delta}_K}}
\]
where $\boldsymbol{\delta}_K \in \mathbb{R}^n$ is derived from the Kyber768
shared secret $K$ via SHAKE-256 (SHA3-256 iterated):
\[
\boldsymbol{\delta}_K[i] = \frac{\text{SHAKE256}(K \| \text{hostname} \| i)}{2^{64}} - 0.5
\]
\end{mathbox}

Since $K$ is a Kyber768 shared secret, computing $\boldsymbol{\delta}_K$
requires solving Module LWE — which requires exponential time classically
and is believed to be quantum-resistant.

\textbf{Effect:} Two clients running identical workloads on identical hardware
have different ``normal'' fidelity values because $\boldsymbol{\delta}_K$
is different for each. An attacker who extracts the algorithm from one
client cannot use it to predict fidelity values for another client.

% ============================================================
\section{Measurement, PDF Generation, and the Complete Pipeline}
% ============================================================

\subsection{What TAARA Measures and How}

\subsubsection{Organizational Health Analysis (OHA) — Pillar A}

The OHA scan connects to the target via SSH (or cloud API) and collects:

\begin{enumerate}
  \item \textbf{sshd\_config} — parsed for all misconfigurations.
  \item \textbf{Firewall state} — \texttt{ufw status}, \texttt{iptables -L}.
  \item \textbf{Open ports} — \texttt{ss -tulnp}.
  \item \textbf{Auth log features} — failed/accepted logins, unique IPs.
  \item \textbf{User accounts} — UID 0 accounts, empty passwords.
\end{enumerate}

Each finding is classified by severity:
\[
\text{Security Risk Score} = \min\left(25 \cdot n_{\text{critical}} + 15 \cdot n_{\text{high}} + 5 \cdot n_{\text{medium}} + n_{\text{low}},\; 100\right)
\]

\subsubsection{Infrastructure Health Score}

The final score combining all three pillars:
\[
\text{IHS} = \max\left(0,\; 100 - \lfloor 0.40 \cdot S_{\text{attack}} + 0.20 \cdot N_{\text{novelty}} + 0.25 \cdot R_{\text{repo}} + 0.15 \cdot C_{\text{cost}} \rfloor\right)
\]

where:
\begin{itemize}
  \item $S_{\text{attack}}$: Security Risk Score from Pillar A
  \item $N_{\text{novelty}}$: Behavior Novelty Score (= $100 \times (1 - F_{\min})$)
  \item $R_{\text{repo}}$: Repository risk score from Pillar B
  \item $C_{\text{cost}}$: Cloud waste score (savings / 500 × 15, capped at 15)
\end{itemize}

\subsubsection{Repository Risk (Pillar B) — The 11-Feature Posture Vector}

For code scanning, TAARA builds a posture vector $\mathbf{p} \in \mathbb{R}^{11}$,
normalized to unit length:

\begin{align*}
p_0 &= \min(n_{\text{direct}} / 100,\; 1.0) & &\text{direct dependencies} \\
p_1 &= \min(n_{\text{transitive}} / 1000,\; 1.0) & &\text{transitive deps} \\
p_2 &= \min(n_{\text{critical}} / 50,\; 1.0) & &\text{critical CVEs} \\
p_3 &= \min(n_{\text{high}} / 100,\; 1.0) & &\text{high CVEs} \\
p_4 &= \min(n_{\text{medium}} / 200,\; 1.0) & &\text{medium CVEs} \\
p_5 &= \min(n_{\text{low}} / 200,\; 1.0) & &\text{low CVEs} \\
p_6 &= n_{\text{affected}} / n_{\text{total}} & &\text{fraction with CVEs} \\
p_7 &= \min(\bar{d} / 5,\; 1.0) & &\text{mean dependency depth} \\
p_8 &= \min(\rho \cdot 100,\; 1.0) & &\text{graph density} \\
p_9 &= \min(\text{chain\_score} / 10,\; 1.0) & &\text{max exploit chain} \\
p_{10} &= 0 & &\text{EOL package count (reserved)}
\end{align*}

Quantum fidelity $F(\mathbf{p}, \mathbf{p}_{\text{policy}})$ tells TAARA how
close the repo's posture is to the ``secure baseline'' captured by the
nearest policy embedding.

\subsubsection{Exploit Chain Scoring}

For each vulnerable package, TAARA scores the exploitability as:
\[
\text{chain\_score} = s_{\text{severity}} \times \frac{1}{d}
\]
where $s \in \{10.0, 7.5, 5.0, 2.5\}$ for \{critical, high, medium, low\}
and $d$ is the dependency depth (1 = direct, 2 = transitive, etc.).

\textbf{Intuition:} A critical CVE in a direct dependency (depth=1) scores
$10.0 \times 1.0 = 10.0$ (maximum urgency). The same CVE in a 4th-level
transitive dependency scores $10.0 \times 0.25 = 2.5$ — still critical,
but harder to reach.

\begin{examplebox}{NodeGoat Scan Results (Real Benchmark Data)}
NodeGoat (OWASP vulnerable Node.js app) scanned with TAARA:
\begin{itemize}
  \item Packages resolved: $\approx 800$ (from \texttt{package-lock.json})
  \item CVEs found: 165 findings, 46 critical, 63 high
  \item Repo quantum fidelity $F = 0.4277 < 0.5$: \textbf{UNSAFE}
  \item Packages flagged include: \texttt{growl 1.9.2} (GHSA-7g98-r9q6-p2wx),
        \texttt{set-value 0.4.3}, \texttt{handlebars 4.0.6}
  \item GraphRAG LLM analysis (Groq LLaMA-3.3-70B): identified these specific
        packages as the highest-risk attack paths, with npm upgrade commands
\end{itemize}
\end{examplebox}

\subsection{TaaraWords PDF Report Generation}

The PDF is generated by ReportLab, a pure-Python PDF library. The pipeline:

\begin{enumerate}
  \item \textbf{Data extraction:} Read directly from \texttt{repo\_results}
        (OSV CVEs, exploit chains, quantum fidelity, cross-file chains).
  \item \textbf{Groq-generated summary:} Call LLaMA-3.3-70B with the real
        CVE findings as context. Strictly prompt-engineered: ``Use only
        the findings provided below. Do not invent vulnerabilities.''
        The LLM converts findings to business-readable language.
  \item \textbf{Quantum fidelity on the summary:} Embed the generated summary,
        compare to policy space. The PDF includes the fidelity score for the
        executive summary itself — a self-auditing report.
  \item \textbf{PDF assembly:} Cover page (risk score + F score), executive
        summary, findings table (top 20 critical, top 15 high), action plan
        (WEEK 1 / MONTH 1 / QUARTER), TaaraWare pitch, data sources.
\end{enumerate}

\textbf{Why HTML escaping matters:} CVE descriptions contain characters like
\texttt{<}, \texttt{>}, \texttt{\&} (e.g., ``XSS via \texttt{<script>}'').
ReportLab uses an XML-like markup language (\texttt{Paragraph}).
Without \texttt{html.escape()}, one CVE description containing ``\texttt{<}''
breaks the entire PDF render with
``\texttt{paraparser: syntax error: parse ended with 1 unclosed tags}''.

\subsection{The Complete System Flow}

\begin{figure}[H]
\centering
\begin{tikzpicture}[
  scale=0.85,
  every node/.style={font=\footnotesize},
  box/.style={draw, rounded corners, minimum width=3cm, minimum height=0.7cm,
               fill=taaraBlue!10, draw=taaraBlue},
  arrow/.style={->, >=Stealth, thick, taaraBlue}
]
  % Input
  \node[box] (ssh) at (0,0) {SSH Server};
  \node[box] (repo) at (0,-1.5) {Git Repo};
  \node[box] (cloud) at (0,-3.0) {Cloud (AWS/GCP)};

  % Collectors
  \node[box, fill=taaraYellow!20] (plat) at (4,0) {Platform Scan};
  \node[box, fill=taaraYellow!20] (scanr) at (4,-1.5) {Repo Scanner};
  \node[box, fill=taaraYellow!20] (cost) at (4,-3.0) {Cost Analyzer};

  % Analysis
  \node[box, fill=taaraRed!15, draw=taaraRed] (ae) at (8,0) {Autoencoder\\+ Memory};
  \node[box, fill=taaraRed!15, draw=taaraRed] (qe) at (8,-1.0) {Quantum Circuit\\$F_{\min}(t)$};
  \node[box, fill=taaraRed!15, draw=taaraRed] (kb) at (8,-2.0) {FAISS + GraphRAG\\+ LLM};
  \node[box, fill=taaraRed!15, draw=taaraRed] (dep) at (8,-3.0) {Dep. Graph\\+ OSV};

  % Scores
  \node[box, fill=taaraGreen!15, draw=taaraGreen!60!black] (ihs) at (12,-1.5)
    {IHS Score\\+ PDF Report};

  % Arrows
  \draw[arrow] (ssh) -- (plat);
  \draw[arrow] (repo) -- (scanr);
  \draw[arrow] (cloud) -- (cost);
  \draw[arrow] (plat) -- (ae);
  \draw[arrow] (plat) -- (kb);
  \draw[arrow] (scanr) -- (dep);
  \draw[arrow] (scanr) -- (kb);
  \draw[arrow] (ae) -- (qe);
  \draw[arrow] (qe) -- (ihs);
  \draw[arrow] (kb) -- (ihs);
  \draw[arrow] (dep) -- (ihs);
  \draw[arrow] (cost) -- (ihs);
\end{tikzpicture}
\caption{Complete TAARA Q.0 data flow from infrastructure to IHS score and PDF.}
\end{figure}

% ============================================================
\section*{Appendix A: Key Code References}
% ============================================================
\addcontentsline{toc}{section}{Appendix A: Key Code References}

\begin{table}[H]
\centering
\footnotesize
\begin{tabular}{@{}llp{7cm}@{}}
\toprule
File & Lines & What is here \\
\midrule
\texttt{components/taara\_core.py} & 59--83 & Least-squares projection (Equations 3–5) \\
\texttt{components/taara\_core.py} & 85--129 & Threshold-free novelty criterion (Eq. 6) \\
\texttt{components/quantum\_engine.py} & 65--93 & 4-qubit PennyLane circuit \\
\texttt{components/quantum\_engine.py} & 135--144 & Fidelity computation $F = |\braket{\psi_t}{\psi_m}|^2$ \\
\texttt{components/quantum\_engine.py} & 353--385 & Kyber-keyed client transform \\
\texttt{components/dna\_autoencoder.py} & 31--67 & Autoencoder architecture \\
\texttt{components/dna\_autoencoder.py} & 86--173 & Training loop with early stopping \\
\texttt{components/taaraware\_manager.py} & 252--374 & Three hidden features \\
\texttt{components/taaraware\_manager.py} & 586--618 & Kyber768 key generation + deployment \\
\texttt{research/query\_knowledge\_base.py} & 46--69 & Fidelity circuit for GraphRAG \\
\texttt{research/query\_knowledge\_base.py} & 134--371 & TAARAScan: full GraphRAG pipeline \\
\texttt{research/scan\_repo.py} & 451--590 & Dependency graph + exploit chains \\
\texttt{research/scan\_repo.py} & 1294--1450 & GraphRAG + LLM + quantum in repo scanner \\
\texttt{components/taara\_analysis.py} & 25--52 & MITRE ATT\&CK mapping table \\
\bottomrule
\end{tabular}
\caption{Where to find each technical component in the codebase.}
\end{table}

% ============================================================
\section*{Appendix B: Glossary}
% ============================================================
\addcontentsline{toc}{section}{Appendix B: Glossary}

\begin{description}
  \item[Amplitude encoding] Encoding a classical vector $\mathbf{v}$ into
    a quantum state by using its normalized components as probability amplitudes.
  \item[CVSS] Common Vulnerability Scoring System. Scores CVEs from 0–10.
    TAARA maps: ${\geq}9.0\to$critical, ${\geq}7.0\to$high, ${\geq}4.0\to$medium.
  \item[FAISS] Facebook AI Similarity Search. Indexes high-dimensional vectors for
    fast approximate nearest-neighbor retrieval.
  \item[Fidelity] $F(\ket{\psi_a}, \ket{\psi_b}) = |\braket{\psi_a}{\psi_b}|^2$.
    1.0 = identical, 0.0 = orthogonal. TAARA uses 0.5 as the safe/unsafe threshold.
  \item[GraphRAG] RAG with a knowledge graph providing causal, not just semantic, context.
  \item[Kyber768 / ML-KEM] NIST FIPS 203 Level 3 post-quantum key encapsulation.
    Quantum-resistant; security based on Module LWE hardness.
  \item[Memory basis] $\mathcal{M}_u$: the set of prior behavioral state vectors
    for identity $u$. Used for reconstruction in §3.
  \item[M-LWE] Module Learning With Errors. The mathematical problem underlying
    Kyber768. Believed hard for both classical and quantum computers.
  \item[OSV.dev] Open Source Vulnerabilities database. TAARA queries it for real CVEs.
  \item[Posture vector] The 11-feature normalized vector summarizing a repo's
    dependency security posture, used for repo-level quantum fidelity scoring.
  \item[Reconstruction residual] $\boldsymbol{\Delta}_t = \mathbf{x}_t - \hat{\mathbf{x}}_t$.
    The component of current behavior not representable by prior observations.
  \item[TaaraWare] Always-on collector agent deployed to customer infrastructure.
    Collects behavioral features and sends to Command Center.
\end{description}

\end{document}
